%!LW recipe=pdflatex x3
\documentclass[10pt,a4paper]{article}

\usepackage[spanish,es-nodecimaldot]{babel}
\usepackage[utf8]{inputenc}
\usepackage[T1]{fontenc}
\usepackage{newtxtext,newtxmath}
\usepackage{geometry}
\usepackage{microtype}
\usepackage{amsmath,mathtools}
\usepackage{enumitem}
\usepackage{xcolor}
\usepackage{hyperref}
\usepackage{fancyhdr}
\usepackage{lastpage}
\usepackage[most]{tcolorbox}

\geometry{margin=2.5cm,includeheadfoot}
\setlength{\headheight}{15pt}
\setlength{\headsep}{0.28cm}
\setlength{\footskip}{0.62cm}
\setlength{\parindent}{0pt}
\setlength{\parskip}{0.28em}
\setlength{\abovedisplayskip}{0.5em}
\setlength{\belowdisplayskip}{0.5em}
\setlength{\abovedisplayshortskip}{0.25em}
\setlength{\belowdisplayshortskip}{0.25em}
\setlist[enumerate]{leftmargin=2.2em,itemsep=0.3em,topsep=0.25em}
\setlist[itemize]{leftmargin=1.5em,itemsep=0.2em,topsep=0.2em}

\definecolor{repblue}{gray}{0}
\definecolor{repteal}{gray}{0}
\definecolor{repgray}{gray}{0}
\definecolor{repaccent}{gray}{0}
\definecolor{repline}{gray}{0}
\definecolor{replight}{gray}{1}

\hypersetup{
  hidelinks,
  pdftitle={TRGF Padilla Tarea 7},
  pdfauthor={Juan Ignacio Padilla Barrientos}
}

\pagestyle{fancy}
\fancyhf{}
\fancyhead[L]{\textcolor{repblue}{\small\itshape TRGF Padilla}}
\fancyhead[R]{\textcolor{repgray}{\small Seminario UCR I 2026}}
\fancyfoot[C]{\textcolor{repgray}{\small \thepage/\pageref*{LastPage}}}
\renewcommand{\headrulewidth}{0.3pt}
\renewcommand{\footrulewidth}{0pt}

\newcommand{\CC}{\mathbb{C}}
\newcommand{\RR}{\mathbb{R}}

\tcbset{
  enhanced,
  sharp corners,
  boxrule=0.45pt,
  arc=0pt,
  left=0.85em,
  right=0.85em,
  top=0.45em,
  bottom=0.4em,
  boxsep=0pt
}

\newcommand{\inlinehead}[2][repblue]{%
  \par\smallskip
  {\color{#1}\bfseries #2.}\hspace{0.45em}%
}

\newcommand{\answer}{\inlinehead{Respuesta}}

\newcommand{\topbanner}{%
  \begin{tcolorbox}[
    colback=replight,
    colframe=repline,
    borderline west={1.2pt}{0pt}{repblue},
    borderline east={1.2pt}{0pt}{repteal},
    left=1.0em,
    right=1.0em,
    top=0.7em,
    bottom=0.55em
  ]
    {\Large\bfseries\color{repblue} TRGF Padilla \textemdash\ Tarea 7}\hfill
    {\color{repaccent}\small\itshape Problemas}\\[-0.1em]
    {\normalsize\color{repteal} Seminario UCR I 2026}
  \end{tcolorbox}
}

\newcounter{probcnt}

\newcommand{\problem}[2]{%
  \stepcounter{probcnt}%
  \ifnum\value{probcnt}>1\newpage\fi
  \par\vspace{0.25em}
  {\large\bfseries\color{repblue} Problema #1}\par
  {\color{repline}\rule{\linewidth}{0.55pt}}\par
  \begin{tcolorbox}[
    breakable,
    colback=white,
    colframe=repline,
    borderline west={1.6pt}{0pt}{repteal},
    left=0.7em,
    right=0.7em,
    top=0.35em,
    bottom=0.3em,
    before skip=0.35em,
    after skip=0.45em
  ]
    #2
  \end{tcolorbox}
}

\begin{document}

\color{black}
\topbanner
\small

\problem{1}{%
Sean $(U,\pi)$ y $(V,\phi)$ dos irreps de un grupo finito $G$. Sea
$L_0:U\to V$ cualquier transformacion lineal. Demuestre que la transformacion
$L:U\to V$ dada por
\[
  L(u):=\frac{1}{|G|}\sum_{g\in G}\phi(g^{-1})L_0\bigl(\pi(g)u\bigr)
\]
es igual a $0$ si $U\not\cong V$ y vale
\[
  \frac{\operatorname{Tr}(L_0)}{\dim U}\operatorname{Id}_U
\]
si $U=V$.
}

\answer
Primero observemos que $L$ entrelaza las acciones de $G$. En efecto, si
$h\in G$ y $u\in U$, entonces
\[
  L\bigl(\pi(h)u\bigr)
  =\frac{1}{|G|}\sum_{g\in G}\phi(g^{-1})L_0\bigl(\pi(gh)u\bigr).
\]
Haciendo el cambio de variable $k=gh$, que recorre todo $G$ cuando $g$ recorre
todo $G$, obtenemos
\[
  L\bigl(\pi(h)u\bigr)
  =\frac{1}{|G|}\sum_{k\in G}\phi\bigl((kh^{-1})^{-1}\bigr)L_0\bigl(\pi(k)u\bigr)
  =\frac{1}{|G|}\sum_{k\in G}\phi(hk^{-1})L_0\bigl(\pi(k)u\bigr).
\]
Como $\phi$ es representación,
\[
  L\bigl(\pi(h)u\bigr)
  =\phi(h)\frac{1}{|G|}\sum_{k\in G}\phi(k^{-1})L_0\bigl(\pi(k)u\bigr)
  =\phi(h)L(u).
\]
Por tanto,
\[
  L\circ\pi(h)=\phi(h)\circ L
  \qquad (h\in G),
\]
es decir, $L\in \operatorname{Hom}_G(U,V)$.

Si $U\not\cong V$, el Lema de Schur implica que
\[
  \operatorname{Hom}_G(U,V)=0,
\]
y por consiguiente $L=0$.

Supongamos ahora que $U=V$; en este caso $L_0\in \operatorname{End}(U)$ y la
fórmula define un endomorfismo $G$-equivariante de $U$. Por el Lema de Schur,
existe $\lambda\in\CC$ tal que
\[
  L=\lambda\operatorname{Id}_U.
\]
Tomando trazas en la fórmula de $L$,
\[
  \operatorname{Tr}(L)
  =\frac{1}{|G|}\sum_{g\in G}\operatorname{Tr}\bigl(\pi(g^{-1})L_0\pi(g)\bigr).
\]
Usando la invariancia cíclica de la traza,
\[
  \operatorname{Tr}\bigl(\pi(g^{-1})L_0\pi(g)\bigr)
  =\operatorname{Tr}\bigl(L_0\pi(g)\pi(g^{-1})\bigr)
  =\operatorname{Tr}(L_0)
\]
para todo $g\in G$. Así,
\[
  \operatorname{Tr}(L)
  =\frac{1}{|G|}\sum_{g\in G}\operatorname{Tr}(L_0)
  =\operatorname{Tr}(L_0).
\]
Por otra parte, como $L=\lambda\operatorname{Id}_U$,
\[
  \operatorname{Tr}(L)=\lambda\dim U.
\]
Luego
\[
  \lambda=\frac{\operatorname{Tr}(L_0)}{\dim U}.
\]
Concluimos que, cuando $U=V$,
\[
  L=\frac{\operatorname{Tr}(L_0)}{\dim U}\operatorname{Id}_U.
\]
Esto prueba ambos casos del enunciado.\hfill$\blacksquare$

\problem{2}{%
Sea $(U,\pi)$ una representacion irreducible de un grupo finito $G$.
Demuestre que el homomorfismo de algebras
\[
  \pi:\CC[G]\to \operatorname{End}(U)
\]
es sobreyectivo.
}

\answer
Sea $n:=\dim U$ y fijemos una base $e_1,\dots,e_n$ de $U$. Para cada
$g\in G$, escribamos
\[
  \pi(g)=\bigl(a_{rs}(g)\bigr)_{1\le r,s\le n}.
\]
Ademas, para $1\le i,j\le n$, sea
\[
  E_{ij}\in \operatorname{End}(U)
\]
la matriz unidad definida por
\[
  E_{ij}(e_j)=e_i,
  \qquad
  E_{ij}(e_m)=0 \quad (m\neq j).
\]

Aplicamos ahora el Problema~1 con $V=U$, $\phi=\pi$ y $L_0=E_{ij}$. Como
\[
  \operatorname{Tr}(E_{ij})=\delta_{ij},
\]
obtenemos
\[
  \frac{1}{|G|}\sum_{g\in G}\pi(g^{-1})E_{ij}\pi(g)
  =\frac{\delta_{ij}}{n}\operatorname{Id}_U.
\]
Tomando la entrada $(r,s)$ de ambos lados, se sigue que
\[
  \frac{1}{|G|}\sum_{g\in G} a_{ri}(g^{-1})\,a_{js}(g)
  =\frac{\delta_{ij}\delta_{rs}}{n}.
\]
Haciendo el cambio de variable $g\mapsto g^{-1}$,
\[
  \frac{1}{|G|}\sum_{g\in G} a_{ri}(g)\,a_{js}(g^{-1})
  =\frac{\delta_{ij}\delta_{rs}}{n}.
\]

Para $1\le p,q\le n$, definamos el elemento
\[
  x_{pq}:=\frac{n}{|G|}\sum_{g\in G} a_{qp}(g^{-1})\,g \in \CC[G].
\]
Entonces
\[
  \pi(x_{pq})
  =\frac{n}{|G|}\sum_{g\in G} a_{qp}(g^{-1})\,\pi(g).
\]
Si tomamos la entrada $(r,s)$ de esta matriz, obtenemos
\[
  \bigl(\pi(x_{pq})\bigr)_{rs}
  =\frac{n}{|G|}\sum_{g\in G} a_{qp}(g^{-1})\,a_{rs}(g).
\]
Aplicando la identidad anterior con $(r,i,j,s)=(r,s,q,p)$, resulta
\[
  \bigl(\pi(x_{pq})\bigr)_{rs}
  =\delta_{rp}\delta_{sq}.
\]
Por tanto,
\[
  \pi(x_{pq})=E_{pq}.
\]

Hemos visto que toda matriz unidad $E_{pq}$ pertenece a la imagen de
\[
  \pi:\CC[G]\to \operatorname{End}(U).
\]
Como las matrices $E_{pq}$ forman una base de $\operatorname{End}(U)$,
concluimos que la imagen de $\pi$ es todo $\operatorname{End}(U)$. En otras
palabras, $\pi$ es sobreyectivo.\hfill$\blacksquare$

\problem{3}{%
Sea $A=\CC[x_g:g\in G]$ el anillo de polinomios que tiene una variable
indeterminada por cada elemento de $G$.

Sea $(U,\pi)$ una representacion irreducible de un grupo finito $G$, fije una
base de $U$ y escriba $n:=\dim U$. En esa base, defina $f_{ij}\in A$, para
$1\le i,j\le n$, como la funcion lineal dada por la entrada $(i,j)$ de la matriz
\[
  \pi\!\left(\sum_{g\in G} x_g g\right).
\]
Demuestre que los polinomios $f_{ij}$ son linealmente independientes en $A$.

\smallskip
\textit{Sugerencia: utilice el problema 1.}
}

\answer
Sea $e_1,\dots,e_n$ la base fija de $U$, y para cada $g\in G$ escribamos
\[
  \pi(g)=\bigl(\pi(g)_{ij}\bigr)_{1\le i,j\le n}.
\]
Como
\[
  \pi\!\left(\sum_{g\in G} x_g g\right)=\sum_{g\in G} x_g\pi(g),
\]
la entrada $(i,j)$ de esa matriz es
\[
  f_{ij}=\sum_{g\in G} x_g\,\pi(g)_{ij}.
\]

Supongamos que existe una relacion lineal
\[
  \sum_{i,j=1}^n \lambda_{ij}f_{ij}=0
  \qquad (\lambda_{ij}\in\CC).
\]
Sustituyendo la expresion anterior para $f_{ij}$, obtenemos
\[
  \sum_{g\in G} x_g\left(\sum_{i,j=1}^n \lambda_{ij}\pi(g)_{ij}\right)=0.
\]
Como las variables $x_g$ son algebraicamente independientes, el coeficiente de
cada $x_g$ debe ser nulo; es decir,
\[
  \sum_{i,j=1}^n \lambda_{ij}\pi(g)_{ij}=0
  \qquad (g\in G).
\]

Para $1\le a,b\le n$, sea $E_{ab}\in\operatorname{End}(U)$ la unidad
matricial dada por
\[
  E_{ab}(e_b)=e_a,
  \qquad
  E_{ab}(e_t)=0\quad (t\ne b).
\]
Aplicando el Problema 1 al caso $U=V$ y $L_0=E_{ab}$, se obtiene
\[
  \frac{1}{|G|}\sum_{g\in G}\pi(g^{-1})E_{ab}\pi(g)
  =\frac{\operatorname{Tr}(E_{ab})}{n}\operatorname{Id}_U
  =\frac{\delta_{ab}}{n}\operatorname{Id}_U.
\]
Tomando la entrada $(k,\ell)$ de ambos lados, resulta
\[
  \frac{1}{|G|}\sum_{g\in G}\pi(g^{-1})_{ka}\,\pi(g)_{b\ell}
  =\frac{\delta_{ab}\delta_{k\ell}}{n}.
\]

Fijemos ahora $a,k\in\{1,\dots,n\}$. Multiplicando la igualdad
\(
\sum_{i,j} \lambda_{ij}\pi(g)_{ij}=0
\)
por $\pi(g^{-1})_{ka}$ y promediando sobre $G$, obtenemos
\[
  0
  =\sum_{i,j=1}^n \lambda_{ij}
    \left(\frac{1}{|G|}\sum_{g\in G}\pi(g^{-1})_{ka}\,\pi(g)_{ij}\right).
\]
Usando la identidad anterior con $(b,\ell)=(i,j)$, se sigue que
\[
  0
  =\sum_{i,j=1}^n \lambda_{ij}\frac{\delta_{ai}\delta_{kj}}{n}
  =\frac{\lambda_{ak}}{n}.
\]
Por tanto $\lambda_{ak}=0$ para todo par $(a,k)$, y en consecuencia todas las
constantes $\lambda_{ij}$ son nulas.

Hemos probado que los polinomios $f_{ij}$ son linealmente independientes en
$A$.\hfill$\blacksquare$

\problem{4}{%
Usando la notacion del problema anterior, demuestre que el polinomio
\[
  \det_U(G)=\det([f_{ij}])\in A
\]
es irreducible.

\smallskip
\textit{Este ejercicio termina de resolver las dudas que teniamos sobre el
determinante de grupo: en efecto, los factores asociados a las irreps son todos
irreducibles.}
}

\answer
Sea $A_1\subseteq A$ el subespacio de los polinomios homogeneos de grado $1$.
Por el Problema 3, las formas lineales
\[
  f_{11},f_{12},\dots,f_{nn}
\]
son linealmente independientes en $A_1$. Como $\dim_\CC A_1=|G|$, podemos
completar esta familia a una base de $A_1$:
\[
  f_{11},f_{12},\dots,f_{nn},h_1,\dots,h_m,
  \qquad m:=|G|-n^2.
\]

Consideremos ahora otro anillo de polinomios
\[
  B:=\CC[y_{ij}:1\le i,j\le n,\ z_1,\dots,z_m].
\]
Como las familias anteriores son bases de los espacios de formas lineales de
$B$ y de $A$, existe un unico isomorfismo de $\CC$-algebras
\[
  \Phi:B\longrightarrow A
\]
tal que
\[
  \Phi(y_{ij})=f_{ij},
  \qquad
  \Phi(z_r)=h_r.
\]
Si definimos
\[
  D_n:=\det\bigl([y_{ij}]_{1\le i,j\le n}\bigr),
\]
entonces
\[
  \Phi(D_n)=\det\bigl([f_{ij}]\bigr)=\det_U(G).
\]
Por consiguiente, basta demostrar que $D_n$ es irreducible. Como $D_n$ no
depende de las variables $z_r$, cualquier factorizacion no trivial de $D_n$ en
$B$ daria una factorizacion no trivial en el subanillo
$\CC[y_{ij}:1\le i,j\le n]$. Asi pues, basta probar la irreducibilidad de
$D_n$ en el anillo de polinomios en las variables $y_{ij}$.

Procedemos por induccion sobre $n$. Si $n=1$, entonces
\[
  D_1=y_{11},
\]
que es irreducible.

Supongamos ahora $n\ge 2$ y que $D_{n-1}$ es irreducible. Sea
\[
  F:=\{y_{11},y_{12},\dots,y_{1n}\}
\]
el conjunto de variables de la primera fila. Todo monomio del determinante
$D_n$ contiene exactamente una variable de la primera fila, de modo que $D_n$
tiene grado total $1$ respecto de las variables de $F$. Si
\[
  D_n=P\,Q
\]
con $P,Q\in\CC[y_{ij}]$, entonces
\[
  \deg_F(P)+\deg_F(Q)=1.
\]
Tras intercambiar los factores si hace falta, podemos suponer que
\[
  \deg_F(Q)=0,
\]
es decir, que $Q$ no depende de ninguna variable de la primera fila.

Expandiendo $D_n$ por la primera fila, obtenemos
\[
  D_n=\sum_{j=1}^n (-1)^{1+j} y_{1j}\Delta_j,
\]
donde $\Delta_j$ es el menor obtenido al eliminar la fila $1$ y la columna
$j$. Como $Q$ no depende de las variables $y_{1j}$ y divide a $D_n$, debe
dividir a cada coeficiente $\Delta_j$ de esta expansion. En particular,
$Q\mid\Delta_1$ y $Q\mid\Delta_2$.

Ahora bien, $\Delta_1$ y $\Delta_2$ son determinantes genericos de tamano
$n-1$, luego por hipotesis inductiva ambos son irreducibles. Ademas, no son
asociados: el polinomio $\Delta_1$ no depende de la variable $y_{21}$, mientras
que $\Delta_2$ si depende de ella; de hecho, el monomio
\[
  y_{21}y_{33}y_{44}\cdots y_{nn}
\]
aparece en $\Delta_2$ con coeficiente $\pm1$. Por tanto, los unicos divisores
comunes de $\Delta_1$ y $\Delta_2$ son unidades.

Pero $Q$ divide a ambos, asi que $Q$ tiene que ser una unidad. Esto contradice
que la factorizacion $D_n=P\,Q$ fuera no trivial. Concluimos que $D_n$ es
irreducible, y como $\det_U(G)=\Phi(D_n)$, tambien $\det_U(G)$ es
irreducible.\hfill$\blacksquare$

\problem{5}{%
Sean $c_1,\dots,c_n$ numeros reales positivos y $A=[a_{ij}]$ una matriz
$n\times n$ que cumple que, para cualesquiera $1\le i_1,i_2\le n$,
\[
  \sum_{j=1}^n c_j a_{i_1 j}\,\overline{a_{i_2 j}}=\delta_{i_1=i_2}.
\]

Es decir, las filas de $A$ son ortonormales con respecto al producto interno que
viene con la forma cuadratica asociada a
$c_1x_1^2+\cdots+c_nx_n^2$.

Demuestre que uno tambien tiene, para cualesquiera $1\le j_1,j_2\le n$,
\[
  \sum_{i=1}^n a_{i j_1}\,\overline{a_{i j_2}}
  =\frac{1}{c_{j_1}}\,\delta_{j_1=j_2}.
\]

\smallskip
\textit{En particular, las columnas de $A$ son ortogonales.}
}

\answer
Sea
\[
  C:=\operatorname{diag}(c_1,\dots,c_n).
\]
La hipotesis del problema dice exactamente que
\[
  A C A^*=I_n,
\]
donde $A^*$ denota la traspuesta conjugada de $A$. En efecto, la entrada
$(i_1,i_2)$ de $A C A^*$ es
\[
  \sum_{j=1}^n a_{i_1j}\,c_j\,\overline{a_{i_2j}},
\]
que por hipotesis vale $\delta_{i_1=i_2}$.

Como cada $c_j$ es positivo, la matriz diagonal
\[
  C^{1/2}:=\operatorname{diag}(\sqrt{c_1},\dots,\sqrt{c_n})
\]
esta bien definida. Pongamos
\[
  B:=A C^{1/2}.
\]
Entonces
\[
  BB^*=A C^{1/2}(C^{1/2})A^*=A C A^*=I_n.
\]
Como $B$ es una matriz cuadrada y satisface $BB^*=I_n$, es invertible; de
hecho $B^{-1}=B^*$. Por tanto tambien se tiene
\[
  B^*B=I_n.
\]
Sustituyendo $B=A C^{1/2}$, obtenemos
\[
  C^{1/2}A^*A C^{1/2}=I_n.
\]
Multiplicando a izquierda y derecha por $C^{-1/2}$, concluimos que
\[
  A^*A=C^{-1}=
  \operatorname{diag}\!\left(\frac1{c_1},\dots,\frac1{c_n}\right).
\]

Ahora miramos la entrada $(j_2,j_1)$ de esta igualdad. Se obtiene
\[
  \sum_{i=1}^n \overline{a_{i j_2}}\,a_{i j_1}
  =\frac{1}{c_{j_1}}\,\delta_{j_1=j_2}.
\]
Tomando conjugados en ambos lados, y usando que el lado derecho es real,
llegamos a
\[
  \sum_{i=1}^n a_{i j_1}\,\overline{a_{i j_2}}
  =\frac{1}{c_{j_1}}\,\delta_{j_1=j_2}.
\]
Esto es exactamente lo que se queria demostrar. En particular, si
$j_1\ne j_2$, la suma vale $0$, de modo que las columnas son ortogonales.
\hfill$\blacksquare$

\problem{6}{%
Sea $G$ un grupo finito, y sean $C_1$ y $C_2$ dos clases conjugadas de $G$, con
$g_1\in C_1$ y $g_2\in C_2$. Utilice el problema anterior para demostrar que
\[
  \sum \chi_{U_i}(g_1)\,\overline{\chi_{U_i}(g_2)}
  =\frac{|G|}{|C_1|}\,\delta_{C_1=C_2},
\]
donde la suma recorre todas las posibles irreps $(U_i,\pi_i)$ de $G$.
}

\answer
Sea $K_1,\dots,K_r$ la lista de todas las clases conjugadas de $G$, y para
cada $s$ fijemos un representante $h_s\in K_s$. Tambien enumeremos las irreps
complejas de $G$ como
\[
  (U_1,\pi_1),\dots,(U_r,\pi_r),
\]
lo cual es posible porque el numero de irreps complejas de un grupo finito es
igual al numero de clases conjugadas.

Consideremos la matriz $M=[m_{is}]$ dada por
\[
  m_{is}:=\chi_{U_i}(h_s),
  \qquad 1\le i,s\le r,
\]
y los numeros positivos
\[
  c_s:=\frac{|K_s|}{|G|}.
\]
La ortogonalidad usual de caracteres irreducibles dice que, para cualesquiera
$1\le i_1,i_2\le r$,
\[
  \frac{1}{|G|}\sum_{g\in G}
  \chi_{U_{i_1}}(g)\,\overline{\chi_{U_{i_2}}(g)}
  =\delta_{i_1=i_2}.
\]
Agrupando la suma por clases conjugadas, esto se reescribe como
\[
  \sum_{s=1}^r \frac{|K_s|}{|G|}
  \chi_{U_{i_1}}(h_s)\,\overline{\chi_{U_{i_2}}(h_s)}
  =\delta_{i_1=i_2}.
\]
Es decir,
\[
  \sum_{s=1}^r c_s\,m_{i_1s}\,\overline{m_{i_2s}}
  =\delta_{i_1=i_2}.
\]
Por tanto, la matriz $M$ satisface exactamente la hipotesis del Problema 5.
Aplicando ese resultado, obtenemos que para cualesquiera $1\le s,t\le r$,
\[
  \sum_{i=1}^r m_{is}\,\overline{m_{it}}
  =\frac{1}{c_s}\,\delta_{s=t}
  =\frac{|G|}{|K_s|}\,\delta_{s=t}.
\]

Finalmente, sean $C_1$ y $C_2$ las clases conjugadas dadas en el enunciado. Si
$C_1=K_s$ y $C_2=K_t$, con representantes $g_1=h_s$ y $g_2=h_t$, la igualdad
anterior se convierte en
\[
  \sum_{i=1}^r \chi_{U_i}(g_1)\,\overline{\chi_{U_i}(g_2)}
  =\frac{|G|}{|C_1|}\,\delta_{C_1=C_2}.
\]
Esto es exactamente la formula pedida.\hfill$\blacksquare$

\end{document}
